% Generated draft: requires mathematical and source review.
\documentclass[11pt]{article}
\usepackage[margin=25mm]{geometry}
\usepackage{fontspec}
\setmainfont{Latin Modern Roman}
\usepackage{amsmath,amssymb,mathtools,unicode-math}
\setmathfont{Latin Modern Math}
\usepackage{enumitem,needspace,xurl,hyperref}
\setlength{\parindent}{0pt}
\setlength{\parskip}{5pt}
\setlength{\emergencystretch}{3em}
\newcommand{\N}{\mathbb{N}}
\newcommand{\Z}{\mathbb{Z}}
\newcommand{\Q}{\mathbb{Q}}
\newcommand{\R}{\mathbb{R}}
\newcommand{\C}{\mathbb{C}}
\newcommand{\F}{\mathbb{F}}
\newcommand{\A}{\mathbb{A}}
\newcommand{\PP}{\mathbb{P}}
\newcommand{\E}{\mathbb{E}}
\newcommand{\eps}{\varepsilon}
\newcommand{\dd}{\,\mathrm{d}}
\begin{document}
\section*{UnsolvedMath ID: 30006657}
\textbf{Source problem code:} catalogue entry
\subsection{Definitions and mathematical statement}
Write \(G=\Z/p\Z\). In the usual Smith-theoretic setting, \(p\) is prime; the supplied record does not explicitly state this restriction. A finite \(G\)-CW complex is a \(G\)-space built from finitely many equivariant cells \(G/H\times D^n\), with \(H\leq G\) and trivial action on \(D^n\). Its fixed-point space is
\[
X^G=\{x\in X: gx=x\text{ for every }g\in G\}.
\]
For a specified commutative coefficient ring \(R\), ordinary homology is denoted by \(H_*(X;R)\). One standard meaning of equivariant cohomology is Borel equivariant cohomology,
\[
H_G^*(X;R)=H^*(EG\times_G X;R),
\]
where \(EG\) is a contractible free \(G\)-CW complex. The coefficient ring and this particular meaning of equivariant cohomology are explanatory parameters, not choices fixed by the source.

\textbf{Reconstruction question.} For a finite \(G\)-CW complex \(X\), determine which homological information about \(X^G\) can be reconstructed from the homological and equivariant-cohomological data of \(X\).

The supplied formulation does not specify the coefficients, the precise input data or retained structure, or the required notion of reconstruction. In particular, it does not determine whether the target is integral homology, mod-\(p\) homology, or another invariant, nor whether the input includes only groups, rings, operations, maps, or a more structured object. Thus this is a reconstruction question requiring further specification, rather than a uniquely determined theorem or conjecture. The original excerpt asks about recovery from the homological information of \(X\); the supplied reconstructed statement additionally mentions equivariant-cohomological data.
\subsection{Short English statement}
For a finite cyclic-group CW complex, ask how much homology of its fixed-point space is determined by homological and equivariant-cohomological information about the original space. The source leaves the coefficients, available structure, and meaning of reconstruction unspecified. Its supplied \texttt{partially\_solved} label is a dated curation assessment, not an independently verified conclusion.
\subsection{Sources}
\begin{itemize}
\item Supplied local UnsolvedMath registry record, permanent ID \(30006657\), title \emph{Recovering Fixed-Point Homology in Smith Theory}, source problem code \texttt{catalogue entry}, in the Oberwolfach Reports Open Problems collection. The proposed year is recorded as 2026; no individual proposer is named.
\item \emph{Cohomology of Finite Groups: Interactions and Applications} (2026), Oberwolfach Reports, pp.\ 541--610, DOI \texttt{10.4171/owr/2026/9}. Supplied report link: \url{https://doi.org/10.4171/owr/2026/9}. No problem-specific page is supplied.
\item Supplied literature reference: Subramanian--Burklund, \emph{Smith theory revisited}, linked through the report DOI: \url{https://doi.org/10.4171/owr/2026/9}.
\item Supplied literature reference: Subramanian, \emph{Categorification of Smith Theory}, \url{https://web.math.ku.dk/~jg/students/subramanian-phdthesis.pdf}.
\item Literature assessment dated exactly \texttt{2026-08-22}: the supplied curation labels the problem partially solved and states that Smith localization and an unstable-module description recover substantial fixed-point cohomology from Borel data, that recent structured work recovers the \(p\)-adic homotopy type in the finite CW setting, and that no complete homological reconstruction is known. These are supplied dated curation claims, not independently verified results or a current literature assessment; their hypotheses and time-sensitive status require independent checking.
\item The registry labels the statement \texttt{corrected\_verified} and reports that it was reconstructed and checked against the official Oberwolfach report. This verification claim is supplied provenance; no independent source consultation or mathematical verification is asserted here.
\end{itemize}
\end{document}
